\section{Related Work}
\label{sec:related}

\textbf{Classical versus learned comparisons.} OSCAR~\cite{wong2021oscar} compares a learned-dynamics operational-space controller with operational-space, inverse-kinematics and joint-space baselines on a Franka Panda. Every policy is trained with PPO, so the classical controller is an action space rather than an independent opponent. The Real Robot Challenge~\cite{bauer2021rrc} found classical methods ahead in its 2020 phase and RL ahead in 2021. Aljalbout et al.~\cite{aljalbout2023action} study 13 control spaces with over 250 RL agents, RL only. The closest work in spirit is Kunapuli et al.~\cite{kunapuli2025leveling}: for quadrotor tracking, the gap between an RL tracker and a geometric controller is much smaller than previously published once objective, tuning data and feedforward are equalised; they use no imitation learner. Samuel and Verma~\cite{samuel2026hybrid} compare proportional-integral-derivative control (PID), a linear-quadratic regulator (LQR), model-predictive control (MPC), soft actor-critic (SAC) and a student distilled from MPC by behaviour cloning (BC) with DAgger (interactive relabelling of the student's own states). The arm is a hybrid rigid-pneumatic one, the simulator is the controllers' own model, and the RL baseline stalls. The student does not beat the teacher (92.4 versus 95.0\% at $N{=}2$ on identical goals) and is no more robust to a mass and inertia mismatch. We differ in using a manipulator whose simulated physics is independent of the controllers' own model, pre-specified rules with a deviation table, and a headroom decomposition (Section~\ref{sec:k3}).

\textbf{Distilling MPC into networks.} Several works distil a model-predictive controller into a network. Nubert et al.~\cite{nubert2019safe} approximate a robust MPC law with a network on a robot manipulator, and Tagliabue and How~\cite{tagliabue2023efficient} learn robust multirotor policies from MPC with tube-guided augmentation. Marquez Julbe et al.~\cite{marquez2025diffusion} use diffusion to imitate multi-modal MPC actions on a 7-DoF arm in a kinematic simulation. There the student has a higher success rate but a larger end-of-run error than the MPC, and no model-error test is run. Our settings S3 and S8 (Table~\ref{tab:settings}) put generative students against an exact or privileged teacher with the same budget accounting as the other methods.

\textbf{Residual learning and generative policies.} Residual policy learning~\cite{silver2018residual,johannink2018residual} adds a learned correction to a hand-designed controller; ResiP~\cite{ankile2024resip} refines a frozen Diffusion Policy with RL. Its appendix reports that an MLP base loses success under exploration noise, converges to similar performance on the easier task and collapses on the harder one. Our exploration-scale lesson (Section~\ref{sec:protocol}) is consistent with that sensitivity. Pan et al.~\cite{pan2025noising} argue that generative control policies owe their advantage to iterative computation with supervised intermediate steps rather than to multi-modality, and that a two-step regression policy (MIP) essentially matches flow policies. One-step flow variants~\cite{li2026ofp,wu2026cola} report large speed-ups (author-reported). Our K4 result is consistent with this direction but underpowered (Section~\ref{sec:k4}).
